\section{Semi-regular representations}
\label{sec:characters_semiregular}

\begin{definition}[Semi-regular representation]%
    \label{def:peps_semiregular}
    \lean{Representation.IsSemiRegular,
        Representation.isSemiRegular_leftRegular}
    \leanok
    A representation of $G$ is semi-regular if it contains every irreducible
    representation of $G$, that is, if every irreducible representation
    admits a nonzero intertwining map into it
    \cite[Definition~4.5]{Schuch2010PEPS}. The left-regular representation
    is semi-regular.
\end{definition}

\begin{lemma}[Semi-regularity under representation equivalence]%
    \label{thm:peps_semiregular_equivalence}
    \lean{Representation.IsSemiRegular.of_equiv,
        Representation.Equiv.isSemiRegular_iff}
    \leanok
    \uses{def:peps_semiregular}
    Equivalent representations $\rho$ and $\sigma$ are semi-regular
    simultaneously. Thus the occurrence condition of
    \cite[Definition~4.5]{Schuch2010PEPS} is independent of coordinates.
\end{lemma}

\begin{proof}\leanok
    If $e$ intertwines $\rho$ with $\sigma$ and is invertible, every
    nonzero intertwiner $f$ from an irreducible representation into
    $\rho$ gives a nonzero intertwiner $ef$ into $\sigma$. Applying the
    same argument to $e^{-1}$ proves the converse.
\end{proof}

\begin{lemma}[Trace duality for a semi-regular representation]%
    \label{thm:peps_semiregular_trace}
    \lean{Representation.character_leftRegular,
        Representation.deltaOperator_leftRegular,
        Representation.sum_irreducibleCharacters_mul_of_isSemiRegular,
        Representation.trace_inv_comp_comp_deltaOperator_of_isSemiRegular}
    \leanok
    \uses{def:peps_delta_operator, def:peps_semiregular}
    If $U$ is semi-regular, then $\tr[U_{g^{-1}}U_h\Delta]=\delta_{g,h}$ for
    all $g,h\in G$, with $U_{g^{-1}}=U_g^\dagger$ for a unitary
    representation. For the left-regular representation
    $\Delta=\frac{1}{|G|}\Id$ \cite[Lemma~4.6]{Schuch2010PEPS}.
\end{lemma}

\begin{proof}\leanok
    \uses{thm:peps_delta_expansion, thm:peps_isotypic_projectors}
    The character of the left-regular representation is $|G|\delta_{k,1}$,
    so every character $\chi_i$ has multiplicity $d_i$ in it, and
    $\Delta=\frac{1}{|G|}\sum_iP_i=\frac{1}{|G|}\Id$. Every
    irreducible representation occurs in the left-regular representation,
    and, since $U$ is semi-regular, in $U$; hence $U$ and the left-regular
    representation have the same irreducible characters. Comparing the two
    expressions for $\tr[U_k\Delta]$ in the left-regular representation gives
    $\sum_id_i\chi_i(k)=|G|\delta_{k,1}$, and therefore
    $\tr[U_{g^{-1}}U_h\Delta]=\frac{1}{|G|}\sum_id_i\chi_i(g^{-1}h)
        =\delta_{g,h}$.
\end{proof}

\begin{lemma}[Twirled representation operators]%
    \label{lem:peps_twirled_representation}
    \lean{Representation.rep_mul_charProjector,
        Representation.apply_conj_of_mem_irreducibleCharacters,
        Representation.sum_rep_conj_apply_of_mem,
        Representation.averageMap_linHom_rep,
        Representation.linearIndependent_charProjector,
        Representation.linearIndependent_of_isSemiRegular,
        Representation.linearIndependent_averageMap_conjClasses}
    \leanok
    \uses{thm:peps_isotypic_projectors, thm:peps_group_averaging,
        thm:peps_semiregular_trace}
    The isotypic projectors $P_i$ commute with every $U_g$, are linearly
    independent, and
    \begin{align*}
        \sigma(U_g)=\frac{1}{|G|}\sum_hU_{hgh^{-1}}
        =\sum_{i=1}^{r}\frac{\chi_i(g)}{\chi_i(1)}\,P_i .
    \end{align*}
    If $U$ is semi-regular, then the operators $U_g$ are linearly independent,
    and so are the twirls $\sigma(U_g)$ of one representative $g$ of each
    conjugacy class \cite[proof of Theorem~4.12]{Schuch2010PEPS}.
\end{lemma}

\begin{proof}\leanok
    On an irreducible subrepresentation with character $\chi_i$ the operator
    $\sum_hU_{hgh^{-1}}$ commutes with the representation, so by Schur's lemma
    it is a scalar, and its trace $|G|\chi_i(g)$ fixes the scalar as
    $|G|\chi_i(g)/\chi_i(1)$; the projector $P_i$ acts there as the identity
    and every other projector as zero. Evaluating $\sum_i\mu_iP_i$ on a nonzero
    vector of such a subrepresentation gives $\mu_i$ times the vector. For
    semi-regular $U$, the functional $X\mapsto\tr[U_{g^{-1}}X\Delta]$ takes the
    value $\delta_{g,h}$ on $U_h$ by
    Lemma~\ref{thm:peps_semiregular_trace}; on $\sigma(U_k)$ it is
    $\frac{1}{|G|}$ times the number of $h$ with $hkh^{-1}=g$, which vanishes
    unless $k$ is conjugate to $g$ and is positive for $k=g$.
\end{proof}

\begin{lemma}[Semi-regularity and the group algebra]
    \label{thm:peps_semiregular_group_algebra}
    \lean{Representation.linearIndependent_iff_injective_asAlgebraHom,
        Representation.isSemiRegular_of_linearIndependent,
        Representation.isSemiRegular_iff_linearIndependent,
        Representation.isSemiRegular_iff_injective_asAlgebraHom}
    \leanok
    \uses{def:peps_semiregular}
    Let $\rho$ be a finite-dimensional complex representation of a finite
    group $G$. The following assertions are equivalent:
    $\rho$ contains every irreducible representation; the operators
    $(\rho(g))_{g\in G}$ are linearly independent; and the induced algebra
    homomorphism $\C[G]\to\operatorname{End}(V)$ is injective.
    This is an algebraic characterization of the occurrence condition in
    \cite[Definition~4.5 and Lemma~4.6]{Schuch2010PEPS}.
    More generally, for a representation of a monoid over a commutative
    semiring, independence of its operators is equivalent to injectivity
    of its monoid-algebra homomorphism.
\end{lemma}

\begin{proof}\leanok
    \uses{thm:peps_semiregular_trace, lem:peps_twirled_representation}
    The trace duality gives independence for a semi-regular
    representation. Conversely, if an irreducible character $\chi$ is
    absent, its character projector vanishes:
    \begin{align}
        \frac{\chi(1)}{|G|}\sum_{g\in G}\chi(g^{-1})\rho(g)=0.
        \label{eq:peps_absent_character}
    \end{align}
    Indeed, the projector vanishes on every irreducible summand of $V$.
    Independence in (\ref{eq:peps_absent_character}) would imply
    $\chi(1)=0$, contrary to the positive dimension of the irreducible
    representation. The standard basis of the group algebra identifies
    the last assertion with independence of the group operators.
\end{proof}

\begin{lemma}[Semi-regularity of the contragredient]
    \label{thm:peps_semiregular_dual}
    \lean{Representation.IsSemiRegular.dual,
        Representation.isSemiRegular_dual_iff}
    \leanok
    \uses{def:peps_semiregular}
    A finite-dimensional complex representation $\rho$ of a finite group
    is semi-regular if and only if its contragredient
    $\rho^*(g)=\rho(g^{-1})^{\mathsf T}$ is semi-regular
    \cite[Definition~5.1 and Lemma~5.2]{Schuch2010PEPS}.
\end{lemma}

\begin{proof}\leanok
    \uses{thm:peps_semiregular_group_algebra}
    Inversion permutes the group indices, and transposition is injective
    on endomorphisms of a finite-dimensional vector space. It therefore
    preserves independence of the group operators. Conversely,
    independence of their transposes implies independence of the
    original operators.
\end{proof}

\begin{lemma}[Irreducibility of the contragredient]
    \label{thm:peps_irreducible_dual}
    \lean{Representation.isIrreducible_dual}
    \leanok
    The contragredient of a finite-dimensional irreducible complex
    representation of a group is irreducible. No unitarity or finite-group
    hypothesis is needed.
\end{lemma}

\begin{proof}\leanok
    The annihilator of an invariant dual subspace is invariant in the
    original representation, since the group action on the dual uses
    the inverse element. Irreducibility makes this annihilator zero or
    the whole space. Finite-dimensional annihilator duality then makes
    the original dual subspace the whole space or zero, respectively.
\end{proof}

\begin{lemma}[Dimension bound for distinct irreducibles]
    \label{thm:peps_semiregular_dimension_bound}
    \lean{Representation.sum_finrank_le_of_distinct_irreducibles_of_isSemiRegular}
    \leanok
    \uses{def:peps_semiregular}
    Let $\sigma$ be a finite-dimensional complex semi-regular representation
    of a finite group $G$. For any finite family $(D_i)_{i\in I}$ of
    irreducible representations with distinct characters,
    $\sum_{i\in I}\dim D_i\leq\dim\sigma$.
    This dimension comparison, which underlies the smallest semi-regular
    representation in \cite[Section~7]{Schuch2010PEPS}, requires no unitarity
    assumption on $\sigma$.
\end{lemma}

\begin{proof}\leanok
    Decomposing $\sigma=\bigoplus_{T\in\mathcal A}T$ as a finite internal sum
    of irreducible subrepresentations, semi-regularity ensures that for each
    $i$ some $T_i\in\mathcal A$ has character $\chi_{D_i}$.
    The distinct characters make $i\mapsto T_i$ injective, and equal
    irreducible characters give equivalent representations with
    $\dim T_i=\dim D_i$.
    Taking dimensions in the internal sum therefore gives
    $\sum_i\dim D_i=\sum_i\dim T_i\leq
        \sum_{T\in\mathcal A}\dim T=\dim\sigma$.
\end{proof}

\begin{lemma}[Tensor products of semi-regular representations]%
    \label{thm:peps_semiregular_tensor_product}
    \lean{Representation.isIrreducible_of_finrank_eq_one,
        Representation.IsSemiRegular.exists_mem_invariants_ne_zero,
        Representation.IsSemiRegular.tprod_of_mem_invariants,
        Representation.IsSemiRegular.tprod}
    \leanok
    \uses{def:peps_semiregular}
    If $U$ and $V$ are semi-regular representations of $G$, then so is
    $U\otimes V$ \cite[proof of Lemma~5.2]{Schuch2010PEPS}.
\end{lemma}

\begin{proof}\leanok
    A one-dimensional representation is irreducible, so the semi-regular
    representation $V$ contains the trivial representation and has a nonzero
    invariant vector $v$. If $D$ is irreducible and $f$ is a nonzero
    intertwining map from $D$ to $U$, then $x\mapsto f(x)\otimes v$ is a
    nonzero intertwining map from $D$ to $U\otimes V$.
\end{proof}

\begin{lemma}[The invariant coevaluation vector]%
    \label{thm:characters_coevaluation_invariance}
    \lean{Representation.averageMap_apply_eq_sum,
        Representation.map_dual_coevaluation,
        contractRight_map_coevaluation, contractRight_map_dual}
    \leanok
    Let $\tau$ be a finite-dimensional complex representation of a group on
    $E$, let $(e_i)$ be a basis of $E$ with dual basis $(e_i^*)$, and let
    $\omega=\sum_ie_i\otimes e_i^*$. Then $\omega$ is invariant under
    $\tau(g)\otimes\tau^*(g)$ for every $g$. Write $\langle c\rangle$ for the
    contraction $e\otimes f\mapsto f(e)$ of $c\in E\otimes E^*$. For operators
    $M,N$ on $E$, $\langle(M\otimes\Id)\omega\rangle=\tr M$ and
    $\langle(N\otimes\tau^*(h))c\rangle=\langle(\tau(h^{-1})N\otimes\Id)c\rangle$. For a finite group, the averaging
    projection is $v\mapsto\frac{1}{|G|}\sum_g\tau(g)v$
    \cite[proof of Lemma~5.2]{Schuch2010PEPS}.
\end{lemma}

\begin{proof}\leanok
    Expanding $\tau^*(g)e_i^*$ in the dual basis,
    $(\tau(g)\otimes\tau^*(g))\,\omega
        =\sum_j\big(\sum_ie_i^*(\tau(g^{-1})e_j)\,\tau(g)e_i\big)\otimes e_j^*
        =\sum_je_j\otimes e_j^*$. The contraction of $(M\otimes\Id)\omega$ is
    $\sum_ie_i^*(Me_i)=\tr M$. The contragredient acts by
    $\tau^*(h)f=f\circ\tau(h^{-1})$, which moves it across the pairing.
\end{proof}
